\section{Knowledge Editing and Continual Learning}

Knowledge editing is concerned with changing a model's factual behavior without retraining the entire model from scratch~\cite{wang2024knowledge,yao2023editing}. At a high level, the problem resembles a targeted form of model updating: one wants to revise a particular factual association while preserving the much larger body of knowledge that is not meant to change. This is what distinguishes model editing from ordinary re-training, where broad redistribution of knowledge across parameters is usually acceptable.

When edits arrive over time, the problem becomes closely related to continual learning. Continual learning studies how a system can absorb new information while retaining previously acquired competence. In the context of model editing, the same tension appears in a more localized form. A model must remain plastic enough to absorb a new fact, yet stable enough not to damage prior edits or unrelated knowledge. The standard continual-learning vocabulary of plasticity, stability, and catastrophic forgetting therefore provides useful background for understanding lifelong model editing, even though the actual interventions may be narrower and more targeted than those usually studied in full continual-training settings.

This connection is important because it clarifies that lifelong editing is not only a question of whether a single fact can be changed. It is also a question of whether a long sequence of such changes can be accumulated without collapse. A method that performs well on one isolated intervention may still fail once updates are applied repeatedly and earlier changes begin to interact with later ones. For that reason, the present thesis treats factual editing as a special case of continual model updating rather than as a purely isolated rewriting problem limited to single-edit success.
